% Part: first-order-logic
% Chapter: axiomatic-deduction
% Section: provability-consistency

\documentclass[../../../include/open-logic-section]{subfiles}

\begin{document}

\iftag{FOL}
      {\olfileid{fol}{axd}{prv}}
      {\olfileid{pl}{axd}{prv}}

\olsection{\usetoken{S}{derivability} మరియు అవైరుధ్యం}

ఇప్పుడు !!{derivability} సంబంధపు అనేక లక్షణాలను స్థాపిస్తాం. అవి
స్వతంత్రంగా కూడా ఆసక్తికరమైనవే; అయితే సంపూర్ణతా సిద్ధాంత నిరూపణలో
ప్రతి ఒక్కటీ ఒక పాత్ర పోషిస్తుంది.

\begin{prop}\ollabel{prop:provability-contr}
  $\Gamma \Proves !A$ అయ్యి, $\Gamma \cup \{!A\}$ విరుద్ధమైతే,
  $\Gamma$ విరుద్ధం.
\end{prop}

\begin{proof}
  $\Gamma \cup \{!A\}$ విరుద్ధమైతే, $\Gamma \cup \{!A\}
  \Proves \lfalse$. \olref[ptn]{prop:reflexivity} వల్ల, ప్రతి
  $!B \in \Gamma$కు $\Gamma \Proves !B$. పరికల్పన ప్రకారం
  $\Gamma \Proves !A$ కూడా కనుక, ప్రతి $!B \in \Gamma \cup
  \{!A\}$కు $\Gamma \Proves !B$. \olref[ptn]{prop:transitivity}
  వల్ల, $\Gamma \Proves \lfalse$; అంటే, $\Gamma$ విరుద్ధం.
\end{proof}

\begin{prop}
\ollabel{prop:prov-incons}
$\Gamma \Proves !A$ అయితే, అప్పుడు మాత్రమే $\Gamma \cup
\{\lnot !A\}$ విరుద్ధం.
\end{prop}

\begin{proof}
మొదట $\Gamma \Proves !A$ అని అనుకుందాం. అప్పుడు
\olref[ptn]{prop:monotonicity} వల్ల $\Gamma \cup \{\lnot !A\}
\Proves !A$. \olref[ptn]{prop:reflexivity} వల్ల $\Gamma \cup
\{\lnot !A\} \Proves \lnot !A$. \olref[prp]{ax:lnot2} వల్ల
$\Proves \lnot !A \lif (!A \lif \lfalse)$ కూడా. కాబట్టి
\olref[ded]{prop:mp}ను రెండుసార్లు వర్తింపజేస్తే, $\Gamma \cup
\{\lnot !A\} \Proves \lfalse$.

ఇప్పుడు $\Gamma \cup \{\lnot !A\}$ విరుద్ధమని, అంటే $\Gamma \cup
\{\lnot !A\} \Proves \lfalse$ అని అనుకుందాం. నిగమన సిద్ధాంతం వల్ల
$\Gamma \Proves \lnot !A \lif \lfalse$. \olref[prp]{ax:lfalse2}
వల్ల $\Gamma \Proves (\lnot !A \lif \lfalse) \lif \lnot\lnot !A$;
కాబట్టి \olref[ded]{prop:mp} వల్ల $\Gamma \Proves \lnot\lnot !A$.
$\Gamma \Proves \lnot\lnot !A \lif !A$ కూడా
(\olref[prp]{ax:dne}) కనుక, \olref[ded]{prop:mp}ను మళ్లీ ఉపయోగించి
$\Gamma \Proves !A$ పొందుతాం.
\end{proof}

\begin{prob}
$\Gamma \Proves \lnot !A$ అయితే, అప్పుడు మాత్రమే $\Gamma \cup
\{!A\}$ విరుద్ధమని నిరూపించండి.
\end{prob}

\begin{prop}\ollabel{prop:explicit-inc}
  $\Gamma \Proves !A$, $\lnot !A \in \Gamma$ అయితే, $\Gamma$
  విరుద్ధం.
\end{prop}

\begin{proof}
  స్వావర్తనత్వం వల్ల $\Gamma \Proves \lnot !A$.
  \olref[prp]{ax:lnot2} వల్ల $\Gamma \Proves \lnot !A \lif
  (!A \lif \lfalse)$. \olref[ded]{prop:mp}ను రెండుసార్లు
  వర్తింపజేస్తే $\Gamma \Proves \lfalse$.
  \sourcecorrection{OLTEAXD-006}{ఆంగ్ల మూలం రెండు మోడస్ పోనెన్స్
  అనువర్తనాలకు అవసరమైన $\Gamma \Proves \lnot !A$ పూర్వపక్షాన్ని
  పేర్కొనలేదు. $\lnot !A \in \Gamma$ నుంచి స్వావర్తనత్వం ద్వారా
  వచ్చే ఆ దశను స్పష్టంగా చేర్చాం.}
\end{proof}


\begin{prop}\ollabel{prop:provability-exhaustive}
  $\Gamma \cup \{!A\}$, $\Gamma \cup \{\lnot !A\}$ రెండూ
  విరుద్ధమైతే, $\Gamma$ విరుద్ధం.
\end{prop}

\begin{proof}
సాధన.
\end{proof}

\tagprob{FOL}
\begin{prob}
  \olref[fol][axd][prv]{prop:provability-exhaustive}ను నిరూపించండి.
\end{prob}
\tagendprob

\tagprob{notFOL}
\begin{prob}
  \olref[pl][axd][prv]{prop:provability-exhaustive}ను నిరూపించండి.
\end{prob}
\tagendprob


\end{document}
